\begin{table}[H]
\centering\small
\caption{Salario mínimo real y salario medio real, enero--agosto de 2021--2026}
\label{tab:minimo-media}
\begin{tabular}{lrrr}
\toprule
Año & Salario mínimo real & Salario medio real & Mínimo como \% del salario medio\\
& \multicolumn{2}{c}{(millones de pesos de 2026)} & \\
\midrule
2021 & 1,33 & 2,24 & 59,3 \\
2022 & 1,34 & 2,21 & 60,7 \\
2023 & 1,38 & 2,21 & 62,6 \\
2024 & 1,44 & 2,23 & 64,7 \\
2025 & 1,51 & 2,28 & 65,9 \\
2026 & 1,75 & 2,57 & 68,2 \\
\bottomrule
\end{tabular}
\fuente{Fuentes: normativa salarial \citep{decreto1785,salarioHistorico,decreto2613,decreto2024,decreto2025,decreto1469}, tabulados del CJC con GEIH \citep{geihTabulados} e IPC del DANE \citep{ipcEmpalme2026}. Los montos se expresan a precios promedio de enero a agosto de 2026. El salario medio corresponde a asalariados de 15 años o más con salario válido positivo, ponderado por FEX\_C18. Se incluyen todas las jornadas, asalariados formales e informales y salarios inferiores al mínimo. El porcentaje se calcula antes de redondear los montos.}
\end{table}
